\section{Instrumentation}
\label{sec:instruments}

\subsection{High frequency suite}
\label{sec:instr-highfreq}

Each tower carries a three axis sonic anemometer and an open path infrared gas
analyser mounted on a boom oriented into the prevailing sector, with a
horizontal separation of 0.18 metres and no vertical separation. Both
instruments are sampled at 20 hertz and timestamped by the datalogger against a
satellite disciplined clock, giving a synchronisation error below 2
milliseconds.

The sensor separation matters for the high frequency loss correction of
Section~\ref{sec:proc-spectral}. At the 4 metre saltmarsh towers the separation
is 4.5 percent of the measurement height, which places the associated cospectral
attenuation at 2 to 4 percent of the flux under unstable conditions and 6 to 9
percent under stable conditions. The correction is applied and its magnitude is
reported per record.

\subsection{Slow response suite}
\label{sec:instr-slow}

The meteorological suite records air temperature and relative humidity in an
aspirated shield, barometric pressure, four component net radiation, soil
temperature at three depths, soil moisture at two depths, soil heat flux from
two plates, and tipping bucket precipitation. Table~\ref{tab:instrument-specs}
gives the specifications and the calibration record.

\begin{table}[t]
\centering
\small
\caption{Instrument specifications and calibration status at the end of the
2024 operating year. Drift is the change in the calibration coefficient between
the two most recent calibrations, expressed as a percentage of reading. The
radiometers were the only instruments to exceed their drift tolerance.}
\label{tab:instrument-specs}
\begin{tabular}{llllr}
\toprule
Quantity & Range & Resolution & Interval & Drift (\%) \\
\midrule
Wind vector        & $\pm 40$ m s$^{-1}$ & 0.01 m s$^{-1}$ & 24 months & 0.4 \\
Sonic temperature  & $-40$ to $+60\,\degC$ & 0.01 K & 24 months & 0.2 \\
CO$_2$ density     & 0 to 60 mmol m$^{-3}$ & 0.02 mmol m$^{-3}$ & 6 months & 1.7 \\
H$_2$O density     & 0 to 2000 mmol m$^{-3}$ & 0.5 mmol m$^{-3}$ & 6 months & 2.1 \\
Air temperature    & $-40$ to $+60\,\degC$ & 0.01 K & 12 months & 0.1 \\
Relative humidity  & 0 to 100 \% & 0.1 \% & 12 months & 1.4 \\
Pressure           & 500 to 1100 hPa & 0.1 hPa & 24 months & 0.1 \\
Net radiation      & $-300$ to $+1500$ W m$^{-2}$ & 0.1 W m$^{-2}$ & 12 months & 3.8 \\
Soil heat flux     & $-200$ to $+200$ W m$^{-2}$ & 0.1 W m$^{-2}$ & 24 months & 2.6 \\
Soil moisture      & 0 to 0.55 m$^{3}$ m$^{-3}$ & 0.001 m$^{3}$ m$^{-3}$ & 24 months & 1.1 \\
Precipitation      & 0 to 300 mm h$^{-1}$ & 0.2 mm & 12 months & 0.9 \\
\bottomrule
\end{tabular}
\end{table}

The net radiometer drift of 3.8 percent exceeds the 2 percent tolerance in the
network protocol \citep{ibls2022protocol} and is the largest single contributor to the energy balance
closure deficit reported in Section~\ref{sec:qc-closure}. All four radiometers
were returned for factory calibration in January 2025 and the 2024 radiation
record has been reprocessed against the returned coefficients with a linear
drift model in time.
